\pdfbookmark[1]{Abstract}{Abstract}
\thispagestyle{empty}
\begin{center}
  \rule{\textwidth}{0.4pt}\\[0.9em]
  {\Large\spacedallcaps{Abstract}}\\[0.7em]
  \rule{\textwidth}{0.4pt}
\end{center}
\bigskip

\noindent An embedded program has no system allocator, and its memory manager can be smaller for
that reason. We present the cellblock of MMgr: a pool fixed at compile time, with a persistent tier
that grows up from the base and a temporary tier that grows down from the size. The two boundaries
are the two halves of one atomic word. A take is one fetch-add on that word, which moves its own
boundary and reads the other in one atomic step. A take that lands past the other boundary adds its
bytes to an overrun that the next writer settles. A give is one store and one fetch-or. Writers hold
the cellblock alone behind a drain of the turnstiles of embedded\_types.

The Lean~4 model proves the invariant of the word through every interleaving of a round, that a
refused take is final, and that two words or an undo break it. It proves the layout of both tiers,
the spill, the trim and the owner word. CBMC, symbolic execution of the machine code at
\texttt{-O0} and \texttt{-O2}, and GenMC under five memory models prove the C against the model.
The manager is 778 bytes of Xtensa at \texttt{-Os}. Four axioms are named.
\cleardoublepage
